\documentclass[12pt]{article}
\usepackage[margin=1in]{geometry}
\usepackage{amsmath, amssymb}
\usepackage{graphicx}
\usepackage{hyperref}
\hypersetup{colorlinks=true, linkcolor=blue, urlcolor=blue}

\begin{document}
\section*{Solution}

There are \(2022\) equally spaced points on a circular track \(\gamma\) of circumference \(2022\).
The points are labeled \(A_1, A_2, \ldots, A_{2022}\) in some order, each label used once.
Initially, Bunbun the Bunny begins at \(A_1\).
She hops along \(\gamma\) from \(A_1\) to \(A_2\), then from \(A_2\) to \(A_3\), until she reaches \(A_{2022}\),
after which she hops back to \(A_1\).
When hopping from \(P\) to \(Q\), she always hops along the shorter of the two arcs \(\widehat{PQ}\) of \(\gamma\);
if \(\overline{PQ}\) is a diameter of \(\gamma\), she moves along either semicircle.

To determine the maximal possible sum of the lengths of the \(2022\) arcs which Bunbun traveled, we consider the following:

Label the points around the circle \(P_1\) to \(P_{2022}\) in circular order. Without loss of generality, let \(A_1 = P_1\).

An equality case occurs when the points are labeled as follows:
\(P_1, P_{1012}, P_2, P_{1013}, \ldots, P_{1011}, P_{2022}\), then back to \(P_1\).

Consider the sequence of points \(A_1 = P_1, A_3, \ldots, A_{2021}\).
The sum of the lengths of the \(2022\) arcs is at most the sum of the major arcs
\(\widehat{A_1A_3}, \widehat{A_3A_5}, \ldots, \widehat{A_{2021}A_1}\).

This is \(2022 \cdot 1011\) minus the sum of the minor arcs
\(\widehat{A_1A_3}, \widehat{A_3A_5}, \ldots, \widehat{A_{2021}A_1}\) (denote this sum as \(S\)).
The sum \(S\) is minimized when \(A_1A_3 \ldots A_{2021}\) forms a convex polygon.
If the polygon includes point \(P_{1012}\) or has points on both sides of the diameter \(P_1P_{1012}\),
the sum of arc lengths is \(2022\). Otherwise, it is
\(P_1P_2P_3 \ldots P_{1011}\) or \(P_1P_{2022}P_{2021} \ldots P_{1013}\),
and the sum of arc lengths is \(2020\).

Thus, the maximal possible sum of the lengths of the \(2022\) arcs is:
\[
2022 \cdot 1011 - 2020 = 2042222.
\]

The answer is: \(\boxed{2042222}\).

\end{document}
